\section{Постановка задачи}
\label{sec:task}

Целью практического занятия является построение конечного автомата, распознающего язык, заданный регулярным выражением, с последующей детерминизацией и минимизацией.

По регулярному выражению варианта требуется:

\begin{itemize}
	\item построить диаграмму переходов конечного автомата;
	\item по диаграмме построить таблицу состояний;
	\item проверить, является ли автомат недетерминированным, и записать объяснение;
	\item если автомат недетерминированный \mdash{} преобразовать его в детерминированный;
	\item проверить, является ли построенный детерминированный автомат минимальным, и записать объяснение;
	\item если автомат не минимален \mdash{} минимизировать его;
	\item если выполнялись преобразования, восстановить регулярное выражение по полученному автомату.
\end{itemize}

Вариант~8. Регулярное выражение:
\[
a2(ab)^{*}(c\mid d)^{+}.
\]
В записи вариантов степень указывается числом после символа, поэтому $a2$ означает $a^{2}$ (ровно два вхождения символа~$a$). Рассматриваемый язык~\mdash{} множество цепочек
\[
L = \bigl\{ aa\,(ab)^{n}\,(c\mid d)^{m} \mid n \ge 0,\; m \ge 1 \bigr\}
\]
над алфавитом $\Sigma = \{a, b, c, d\}$.
